% 3.6 Bonus: scalability experiments.
% Evidence: results-verify.md and results-tasks.md (F11-F15, F21-F22).

\subsection{Bonus: Scalability Experiments}
\label{sec:bonus}

The bonus question has two independent scale axes: 500,000 files stress the
build side and the document table, while 400,000,000 distinct words stress the
dictionary and the posting lists. The project implements three mechanisms for
these axes---SPIMI-style block construction, a disk-based B+ tree dictionary,
and the delta-coded v2 index format---and runs the scale experiments that fit
on one machine.

\subsubsection{Scale Boundaries and Measured Range}
\label{sec:scale-range}

The 500,000-document synthetic corpus contains 4,908,747 distinct terms. Its
full in-memory index peaks at 2.00\,GiB. Scaling that figure linearly by term
count to 400,000,000 distinct words would require about 160\,GiB, which a
single machine of this class cannot provide. An in-memory build therefore
cannot answer the second axis directly.

The file axis is measured end to end: statistics, block construction, and the
in-memory build all run on the same 16\,GiB machine, and no address-space or
document-table problem appears at 500,000 files. The distinct-word axis is not
run end to end. The query side is measured up to 50,000,000 B+ tree keys and
extrapolated from the measured trend; the build side is measured to 4.9
million terms and bounded by the $O(\text{budget})$ memory argument. Measured
points and extrapolated points are labelled as such throughout this section.

\subsubsection{Disk B+ Tree Dictionary}
\label{sec:bptree-results}

The query-side problem is not how to search faster but how to avoid loading
the complete term table. The B+ tree sidecar maps each term to a posting
offset and block length, so a lookup touches only a few pages on the root-to-
leaf path. The experiment measures 100,000 to 50,000,000 keys; results are in
Table~\ref{tab:bptree-scale}.

\begin{table}[!htbp]
  \centering
  \caption{Measured scaling of the B+ tree dictionary (4\,KiB pages, bulk
  load, 1,000 random lookups all verified).}
  \label{tab:bptree-scale}
  \small
  \begin{tabular}{r r r r r r}
    \toprule
    Keys & Height & Pages & File bytes & Bytes/key & Page reads/query \\
    \midrule
    100\,000     & 2 & 369     & 1\,515\,520   & 15.1552 & 2.00 \\
    1\,000\,000   & 3 & 3\,688   & 15\,110\,144  & 15.1101 & 3.00 \\
    10\,000\,000  & 3 & 36\,865  & 151\,003\,136 & 15.1003 & 3.00 \\
    30\,000\,000  & 3 & 110\,593 & 452\,993\,024 & 15.0998 & 3.00 \\
    50\,000\,000  & 4 & 184\,322 & 754\,987\,008 & 15.0997 & 4.00 \\
    \bottomrule
  \end{tabular}
  \par\smallskip
  \footnotesize Note: page reads are the actual number of pages read with a
  one-page cache; they equal the tree height.
\end{table}

% figures/bptree.tex -- B+ 树 词典的规模实验：文件体积与查询页读
%
% 数据：reports/data/csv/verify/bptree-bench.csv（本轮 mse bptree-bench，
% 4 KiB 页；100k/1M/10M/30M/50M 键各一次测量）。由
% figures/data/make_figures_data.py 生成 figures/data/bptree-*.dat 与
% bptree-meta.tex（线性拟合、4 亿键外推）。
% 主文档只写 % figures/bptree.tex -- B+ 树 词典的规模实验：文件体积与查询页读
%
% 数据：reports/data/csv/verify/bptree-bench.csv（本轮 mse bptree-bench，
% 4 KiB 页；100k/1M/10M/30M/50M 键各一次测量）。由
% figures/data/make_figures_data.py 生成 figures/data/bptree-*.dat 与
% bptree-meta.tex（线性拟合、4 亿键外推）。
% 主文档只写 % figures/bptree.tex -- B+ 树 词典的规模实验：文件体积与查询页读
%
% 数据：reports/data/csv/verify/bptree-bench.csv（本轮 mse bptree-bench，
% 4 KiB 页；100k/1M/10M/30M/50M 键各一次测量）。由
% figures/data/make_figures_data.py 生成 figures/data/bptree-*.dat 与
% bptree-meta.tex（线性拟合、4 亿键外推）。
% 主文档只写 \input{figures/bptree}。
\input{figures/data/bptree-meta}
\begin{figure}[!htbp]
\centering
\captionsetup[subfigure]{font=scriptsize,skip=2pt}
% ---------------- (a) file size ----------------
\begin{subfigure}{0.49\linewidth}
  \centering
  \begin{tikzpicture}
  \begin{axis}[
    width=7.10cm, height=4.25cm,
    xmode=log, ymode=log,
    xmin=8e4, xmax=1e9, ymin=1e6, ymax=3e10,
    xlabel={number of keys (log scale)}, ylabel={file size (log scale)},
    xtick={1e5,1e6,1e7,1e8}, xticklabels={100k,1M,10M,100M},
    extra x ticks={4e8}, extra x tick labels={400M},
    extra x tick style={tick label style={font=\scriptsize}},
    ytick={1e6,1e7,1e8,1e9,1e10},
    yticklabels={1 MB,10 MB,100 MB,1 GB,10 GB},
    tick label style={font=\scriptsize},
    label style={font=\small},
    grid=major, grid style={black!8},
    legend style={draw=none, fill=none, font=\scriptsize,
                  at={(0.02,0.98)}, anchor=north west},
    legend cell align=left,
  ]
    % 实测点
    \addplot[only marks, mark=*, mark size=1.4pt, color=mseblue]
      table[x=keys, y=file_bytes] {figures/data/bptree-size.dat};
    \addlegendentry{measured (4 KiB pages)}
    % 15.1 B/key 线性拟合（对数坐标下为直线）
    \addplot[dashed, line width=0.8pt, color=msegray]
      table[x=keys, y=fit_bytes] {figures/data/bptree-size-fit.dat};
    \addlegendentry{linear fit (\BptSlope{} B/key)}
    % 4 亿键外推点（空心）
    \addplot[only marks, mark=o,
             mark options={draw=mseblue, fill=white, line width=0.7pt},
             mark size=2.2pt]
      table[x=keys, y=fit_bytes] {figures/data/bptree-extrap.dat};
    \addlegendentry{400M-key extrapolation}
  \end{axis}
  \end{tikzpicture}
  \caption{文件字节 $=$ \BptSlope{} B/键 $\times$ 键数，严格线性}
\end{subfigure}\hfill
% ---------------- (b) page reads ----------------
\begin{subfigure}{0.49\linewidth}
  \centering
  \begin{tikzpicture}
  \begin{axis}[
    width=7.10cm, height=4.25cm,
    xmode=log,
    xmin=8e4, xmax=1e9, ymin=0, ymax=5,
    xlabel={number of keys (log scale)},
    ylabel={page reads per lookup},
    xtick={1e5,1e6,1e7,1e8}, xticklabels={100k,1M,10M,100M},
    extra x ticks={4e8}, extra x tick labels={400M},
    extra x tick style={tick label style={font=\scriptsize}},
    ytick={0,1,2,3,4,5},
    tick label style={font=\scriptsize},
    label style={font=\small},
    grid=major, grid style={black!8},
    legend style={draw=none, fill=none, font=\scriptsize,
                  at={(0.02,0.98)}, anchor=north west},
    legend cell align=left,
  ]
    % 实测页读（点旁标树高）
    \addplot[only marks, mark=*, mark size=1.4pt, color=mseblue]
      table[x=keys, y=page_reads] {figures/data/bptree-reads.dat};
    \addlegendentry{measured (h $=$ tree height)}
    % 4 亿键外推（虚线延到空心点）
    \addplot[dashed, line width=0.8pt, color=msegray]
      table[x=keys, y=page_reads] {figures/data/bptree-reads-fit.dat};
    \addlegendentry{extrapolated}
    \addplot[only marks, mark=o,
             mark options={draw=mseblue, fill=white, line width=0.7pt},
             mark size=2.2pt]
      table[x=keys, y=page_reads] {figures/data/bptree-extrap.dat};
    % 外推读数标注（放在数据点下方的空白处）
    \node[font=\scriptsize, text=msegray, align=right, anchor=south east]
      at (rel axis cs:0.98,0.04)
      {\BptExtrapM{}M keys: height \BptExtrapH{},\\
       \BptExtrapPages{} page reads $=$ \BptExtrapKiB{} KiB};
  \end{axis}
  \end{tikzpicture}
  \caption{每次查询只读树高页：4 亿键也只需 4 页 $=$ 16 KiB}
\end{subfigure}
\caption{B+ 树 词典的规模实验：文件体积严格线性（\BptSlope{} B/键），而每次
随机查询的页读等于树高、随高度对数增长；按实测扇外出推，4 亿个不同词只需
\BptExtrapH{} 层、每次查询 \BptExtrapPages{} 页（\BptExtrapKiB{} KiB），文件约
\BptExtrapGB{} GB。}%
\label{fig:bptree}
\end{figure}
。
% bptree-meta.tex -- generated by figures/data/make_figures_data.py
\newcommand{\BptSlope}{15.1}
\newcommand{\BptExtrapM}{400}
\newcommand{\BptExtrapGB}{6.04}
\newcommand{\BptExtrapH}{4}
\newcommand{\BptExtrapPages}{4}
\newcommand{\BptExtrapKiB}{16}
\newcommand{\BptLeafKeys}{271}
\newcommand{\BptFanout}{372}

\begin{figure}[!htbp]
\centering
\captionsetup[subfigure]{font=scriptsize,skip=2pt}
% ---------------- (a) file size ----------------
\begin{subfigure}{0.49\linewidth}
  \centering
  \begin{tikzpicture}
  \begin{axis}[
    width=7.10cm, height=4.25cm,
    xmode=log, ymode=log,
    xmin=8e4, xmax=1e9, ymin=1e6, ymax=3e10,
    xlabel={number of keys (log scale)}, ylabel={file size (log scale)},
    xtick={1e5,1e6,1e7,1e8}, xticklabels={100k,1M,10M,100M},
    extra x ticks={4e8}, extra x tick labels={400M},
    extra x tick style={tick label style={font=\scriptsize}},
    ytick={1e6,1e7,1e8,1e9,1e10},
    yticklabels={1 MB,10 MB,100 MB,1 GB,10 GB},
    tick label style={font=\scriptsize},
    label style={font=\small},
    grid=major, grid style={black!8},
    legend style={draw=none, fill=none, font=\scriptsize,
                  at={(0.02,0.98)}, anchor=north west},
    legend cell align=left,
  ]
    % 实测点
    \addplot[only marks, mark=*, mark size=1.4pt, color=mseblue]
      table[x=keys, y=file_bytes] {figures/data/bptree-size.dat};
    \addlegendentry{measured (4 KiB pages)}
    % 15.1 B/key 线性拟合（对数坐标下为直线）
    \addplot[dashed, line width=0.8pt, color=msegray]
      table[x=keys, y=fit_bytes] {figures/data/bptree-size-fit.dat};
    \addlegendentry{linear fit (\BptSlope{} B/key)}
    % 4 亿键外推点（空心）
    \addplot[only marks, mark=o,
             mark options={draw=mseblue, fill=white, line width=0.7pt},
             mark size=2.2pt]
      table[x=keys, y=fit_bytes] {figures/data/bptree-extrap.dat};
    \addlegendentry{400M-key extrapolation}
  \end{axis}
  \end{tikzpicture}
  \caption{文件字节 $=$ \BptSlope{} B/键 $\times$ 键数，严格线性}
\end{subfigure}\hfill
% ---------------- (b) page reads ----------------
\begin{subfigure}{0.49\linewidth}
  \centering
  \begin{tikzpicture}
  \begin{axis}[
    width=7.10cm, height=4.25cm,
    xmode=log,
    xmin=8e4, xmax=1e9, ymin=0, ymax=5,
    xlabel={number of keys (log scale)},
    ylabel={page reads per lookup},
    xtick={1e5,1e6,1e7,1e8}, xticklabels={100k,1M,10M,100M},
    extra x ticks={4e8}, extra x tick labels={400M},
    extra x tick style={tick label style={font=\scriptsize}},
    ytick={0,1,2,3,4,5},
    tick label style={font=\scriptsize},
    label style={font=\small},
    grid=major, grid style={black!8},
    legend style={draw=none, fill=none, font=\scriptsize,
                  at={(0.02,0.98)}, anchor=north west},
    legend cell align=left,
  ]
    % 实测页读（点旁标树高）
    \addplot[only marks, mark=*, mark size=1.4pt, color=mseblue]
      table[x=keys, y=page_reads] {figures/data/bptree-reads.dat};
    \addlegendentry{measured (h $=$ tree height)}
    % 4 亿键外推（虚线延到空心点）
    \addplot[dashed, line width=0.8pt, color=msegray]
      table[x=keys, y=page_reads] {figures/data/bptree-reads-fit.dat};
    \addlegendentry{extrapolated}
    \addplot[only marks, mark=o,
             mark options={draw=mseblue, fill=white, line width=0.7pt},
             mark size=2.2pt]
      table[x=keys, y=page_reads] {figures/data/bptree-extrap.dat};
    % 外推读数标注（放在数据点下方的空白处）
    \node[font=\scriptsize, text=msegray, align=right, anchor=south east]
      at (rel axis cs:0.98,0.04)
      {\BptExtrapM{}M keys: height \BptExtrapH{},\\
       \BptExtrapPages{} page reads $=$ \BptExtrapKiB{} KiB};
  \end{axis}
  \end{tikzpicture}
  \caption{每次查询只读树高页：4 亿键也只需 4 页 $=$ 16 KiB}
\end{subfigure}
\caption{B+ 树 词典的规模实验：文件体积严格线性（\BptSlope{} B/键），而每次
随机查询的页读等于树高、随高度对数增长；按实测扇外出推，4 亿个不同词只需
\BptExtrapH{} 层、每次查询 \BptExtrapPages{} 页（\BptExtrapKiB{} KiB），文件约
\BptExtrapGB{} GB。}%
\label{fig:bptree}
\end{figure}
。
% bptree-meta.tex -- generated by figures/data/make_figures_data.py
\newcommand{\BptSlope}{15.1}
\newcommand{\BptExtrapM}{400}
\newcommand{\BptExtrapGB}{6.04}
\newcommand{\BptExtrapH}{4}
\newcommand{\BptExtrapPages}{4}
\newcommand{\BptExtrapKiB}{16}
\newcommand{\BptLeafKeys}{271}
\newcommand{\BptFanout}{372}

\begin{figure}[!htbp]
\centering
\captionsetup[subfigure]{font=scriptsize,skip=2pt}
% ---------------- (a) file size ----------------
\begin{subfigure}{0.49\linewidth}
  \centering
  \begin{tikzpicture}
  \begin{axis}[
    width=7.10cm, height=4.25cm,
    xmode=log, ymode=log,
    xmin=8e4, xmax=1e9, ymin=1e6, ymax=3e10,
    xlabel={number of keys (log scale)}, ylabel={file size (log scale)},
    xtick={1e5,1e6,1e7,1e8}, xticklabels={100k,1M,10M,100M},
    extra x ticks={4e8}, extra x tick labels={400M},
    extra x tick style={tick label style={font=\scriptsize}},
    ytick={1e6,1e7,1e8,1e9,1e10},
    yticklabels={1 MB,10 MB,100 MB,1 GB,10 GB},
    tick label style={font=\scriptsize},
    label style={font=\small},
    grid=major, grid style={black!8},
    legend style={draw=none, fill=none, font=\scriptsize,
                  at={(0.02,0.98)}, anchor=north west},
    legend cell align=left,
  ]
    % 实测点
    \addplot[only marks, mark=*, mark size=1.4pt, color=mseblue]
      table[x=keys, y=file_bytes] {figures/data/bptree-size.dat};
    \addlegendentry{measured (4 KiB pages)}
    % 15.1 B/key 线性拟合（对数坐标下为直线）
    \addplot[dashed, line width=0.8pt, color=msegray]
      table[x=keys, y=fit_bytes] {figures/data/bptree-size-fit.dat};
    \addlegendentry{linear fit (\BptSlope{} B/key)}
    % 4 亿键外推点（空心）
    \addplot[only marks, mark=o,
             mark options={draw=mseblue, fill=white, line width=0.7pt},
             mark size=2.2pt]
      table[x=keys, y=fit_bytes] {figures/data/bptree-extrap.dat};
    \addlegendentry{400M-key extrapolation}
  \end{axis}
  \end{tikzpicture}
  \caption{文件字节 $=$ \BptSlope{} B/键 $\times$ 键数，严格线性}
\end{subfigure}\hfill
% ---------------- (b) page reads ----------------
\begin{subfigure}{0.49\linewidth}
  \centering
  \begin{tikzpicture}
  \begin{axis}[
    width=7.10cm, height=4.25cm,
    xmode=log,
    xmin=8e4, xmax=1e9, ymin=0, ymax=5,
    xlabel={number of keys (log scale)},
    ylabel={page reads per lookup},
    xtick={1e5,1e6,1e7,1e8}, xticklabels={100k,1M,10M,100M},
    extra x ticks={4e8}, extra x tick labels={400M},
    extra x tick style={tick label style={font=\scriptsize}},
    ytick={0,1,2,3,4,5},
    tick label style={font=\scriptsize},
    label style={font=\small},
    grid=major, grid style={black!8},
    legend style={draw=none, fill=none, font=\scriptsize,
                  at={(0.02,0.98)}, anchor=north west},
    legend cell align=left,
  ]
    % 实测页读（点旁标树高）
    \addplot[only marks, mark=*, mark size=1.4pt, color=mseblue]
      table[x=keys, y=page_reads] {figures/data/bptree-reads.dat};
    \addlegendentry{measured (h $=$ tree height)}
    % 4 亿键外推（虚线延到空心点）
    \addplot[dashed, line width=0.8pt, color=msegray]
      table[x=keys, y=page_reads] {figures/data/bptree-reads-fit.dat};
    \addlegendentry{extrapolated}
    \addplot[only marks, mark=o,
             mark options={draw=mseblue, fill=white, line width=0.7pt},
             mark size=2.2pt]
      table[x=keys, y=page_reads] {figures/data/bptree-extrap.dat};
    % 外推读数标注（放在数据点下方的空白处）
    \node[font=\scriptsize, text=msegray, align=right, anchor=south east]
      at (rel axis cs:0.98,0.04)
      {\BptExtrapM{}M keys: height \BptExtrapH{},\\
       \BptExtrapPages{} page reads $=$ \BptExtrapKiB{} KiB};
  \end{axis}
  \end{tikzpicture}
  \caption{每次查询只读树高页：4 亿键也只需 4 页 $=$ 16 KiB}
\end{subfigure}
\caption{B+ 树 词典的规模实验：文件体积严格线性（\BptSlope{} B/键），而每次
随机查询的页读等于树高、随高度对数增长；按实测扇外出推，4 亿个不同词只需
\BptExtrapH{} 层、每次查询 \BptExtrapPages{} 页（\BptExtrapKiB{} KiB），文件约
\BptExtrapGB{} GB。}%
\label{fig:bptree}
\end{figure}


Across the five measured points, bytes per key stays between 15.10 and 15.16,
so file size is almost exactly linear in the key count. The fanout is about
330--370 and a leaf holds about 272 keys on average, so four levels already
cover $10^8$-scale keys. Extrapolating the measured trend to 400,000,000 keys
gives a file of about 6.04\,GB, still height 4, and four page reads per lookup,
about 16\,KiB. Height and page reads are already measured at 50,000,000 keys;
only the file size and page count are extrapolated.

\subsubsection{Block-Based Construction with 500,000 Documents}
\label{sec:block-results}

On the formal Shakespeare corpus, block construction with a 4,096-term budget
produces 17 temporary runs, a peak of 5,129 live terms, and 4,166,556 bytes of
temporary data; its output is byte-identical to the in-memory build.

The 500,000-document experiment scales this path to realistic small-file size.
The synthetic corpus has 500,000 files, 20,000,000 tokens, and 4,908,747
distinct terms. Block construction produces 4,855 runs; scanning takes 5.47\,s,
merging 125.08\,s, for a total wall-clock time of 168.16\,s and a peak memory
of 1.01\,GiB. The in-memory build takes 76.75\,s and peaks at 2.00\,GiB. Both
produce a v1 index of 427,603,247 bytes with the same SHA-256.

\begin{table}[!htbp]
  \centering
  \caption{Build results for 500,000 documents and 20,000,000 tokens.}
  \label{tab:500k}
  \small
  \begin{tabularx}{\linewidth}{X r r r l}
    \toprule
    Build & Index bytes & Build phase & Peak memory & Output \\
    \midrule
    Block, v1 & 427\,603\,247 & scan 5.47 + merge 125.08\,s & 1.01\,GiB & SHA-256 identical \\
    Block, v2 & 284\,560\,009 & scan 5.42 + merge 123.31\,s & 1.46\,GiB & delta-coded \\
    In-memory, v1 & 427\,603\,247 & build 44.93\,s, total 76.75\,s & 2.00\,GiB & baseline \\
    \bottomrule
  \end{tabularx}
  \par\smallskip
  \footnotesize Note: the block-built and in-memory v1 indexes are identical
  under \code{cmp} and SHA-256.
\end{table}

Block construction almost halves peak memory, at the cost of a longer merge.
The current merge finds the smallest term by scanning all 4,855 cursors, about
$O(T\cdot R)$; a heap-based $k$-way merge would reduce this to
$O(T\log R)$. The merge time and the temporary space are the engineering price
of the memory bound.

\subsubsection{Compression Results of the v2 Index Format}
\label{sec:compression}

The v2 format changes only the posting region: document IDs and positions are
delta-coded, and term frequencies and deltas use byte-aligned variable-length
codes. On the main corpus the index drops from 3,685,655 to 1,725,694 bytes,
46.82\% of v1, while build time stays at 0.19\,s and query output is unchanged.

On the 500,000-document corpus the posting region is only 56.1\% of the total
file, so the term table becomes the main contributor to size and the overall
ratio is worse: 427,603,247 bytes fall to 284,560,009 bytes, 66.55\% of v1.
The posting region itself compresses to about 59.6\%, close to the 61.8\% on the
main corpus.

\begin{table}[!htbp]
  \centering
  \caption{v2 compression results and the main cost.}
  \label{tab:compression}
  \small
  \begin{tabularx}{\linewidth}{X r r r r l}
    \toprule
    Corpus & v1 bytes & v2 bytes & v2/v1 & Saving & Cost \\
    \midrule
    Shakespeare & 3\,685\,655 & 1\,725\,694 & 46.82\% & 53.18\% & 0.19\,s build \\
    500k synthetic & 427\,603\,247 & 284\,560\,009 & 66.55\% & 33.45\% & 1.46\,GiB peak (v1: 1.01) \\
    \bottomrule
  \end{tabularx}
  \par\smallskip
  \footnotesize Note: the 500k ratio is dominated by the term table, which v2
  does not compress.
\end{table}

Compression suits storage-constrained settings where CPU time and build memory
are relatively abundant. The 3.7\,MB main index fits in memory entirely, so the
v2 benefit is mainly on disk and in transfer; on the 500k corpus the compressed
build raises peak memory from 1.01 to 1.46\,GiB, so compression does not win on
every resource dimension. These trade-offs are analyzed further in
Section~\ref{sec:bonus-analysis}.
